\documentclass{ximera}
\title{Summation}
\begin{document}
\begin{abstract}
Sigma notation for sums, its basic properties (including the telescoping sum and the formulas for $\sum i$, $\sum i^2$, $\sum i^3$), and some computations with it. Textbook Appendix E.
\end{abstract}
\maketitle

We're going to kind of change topics. Remember that when we're looking at a graph, we used derivatives to find the slope at any point. We're now going to change our question. How can we calculate the \wordChoice{\choice[correct]{area}\choice{slope}\choice{tangent line}\choice{rate of change}}?

Most of the time, it's easy. For a square we have the length times the height. For a triangle we have one half the base times the height. For a circle we have $\pi r^{2}$. But what if we want to calculate the area under a curve?

Before looking fully into this topic, we're going to take a quick detour into summations.

\begin{example}\label{example:12-1}
Say that we want to add 1 to itself ten times. We can do it like: $1+1+1+1+1+1+1+1+1+1=10$. But this becomes more complicated if we want to do more than ten. Say we want to do it 25 times:
\[
1+1+1+1+1+1+1+1+1+1+1+1+1+1+1+1+1+1+1+1+1+1+1+1+1=25 .
\]

This is way to long and almost imposible to read. So we want to rewrite this in a more visible manner. For adding 1 to itself, it's pretty easy: we do $1 \cdot n$ where $n$ ist he number of times we want to add 1 to itself.

How about if we want to add every integer from 1 to 10 : $1+2+3+4+5+$ $6+7+8+9+10=55$. What if we want to do from 1 to 25 ?
% work: 25*26/2 = 325
\[
1+2+3+4+5+6+7+8+9+10+11+12+13+14+15+16+17+18+19+20+21+22+23+24+25=\answer{325}
\]
\end{example}

This is way to complicated and we want to condense this summation by introducing a new notation \wordChoice{\choice[correct]{$\Sigma$ (sigma notation)}\choice{$\Pi$ (product notation)}\choice{$\int$ (integral notation)}\choice{$\Delta$ (difference notation)}}.

\begin{example}\label{example:12-2}
Here are the two examples from before.
\[
\sum_{i=1}^{10} 1=10 \quad \sum_{i=1}^{10} i=55
\]

In general, we have something like the following
\[
\sum_{i=m}^{n} a_{i}=\answer{a_{m}+a_{m+1}+\cdots+a_{n}}
\]
\end{example}

The $i$ are called the indices of summation and they can be thought of as functions.

We can also do this sum without end! If we want to take the summation forever then we let $n=\infty: \sum_{i=m}^{\infty} a_{i}$.

\begin{example}\label{example:12-3}
Let's do some examples.
% work: sum_{i=3}^{5} (i^2+1) = 10 + 17 + 26 = 53;  sum_{j=-4}^{-1} 3 = 3*4 = 12
\[
\begin{aligned}
\sum_{i=m}^{n} a_{i} & =\sum_{i=1}^{4} 2 i-1=2(1)-1+2(2)-1+2(3)-1+2(4)-1=2+4+6+8-4=16 \\
\sum_{i=m}^{n} b_{i} & =\sum_{i=3}^{5} i^{2}+1= \answer{53} \\
\sum_{j=x}^{y} i_{j}=\sum_{j=-4}^{-1} 3 & = \answer{12}
\end{aligned}
\]

And in the other direction, we have
% work: 5,8,11,14 = 3i+2 for i=1..4;  1,0,-1,-2,-3 = 2-i for i=1..5
\[
\begin{gathered}
2+2+2+2+2=\sum_{i=1}^{5} 2 \\
5+8+11+14=\answer{\sum_{i=1}^{4} (3i+2)} \\
1+0-1-2-3=\answer{\sum_{i=1}^{5} (2-i)}
\end{gathered}
\]
\end{example}

Here rae some properties of summations using sigma notation:

\begin{proposition}\label{proposition:12-4}
\begin{enumerate}
  \item $\sum_{i=m}^{n} a_{i}=\sum_{j=m}^{n} a_{j}$.
  \item $\sum_{i=m}^{n} a_{i}=\sum_{i=1}^{n-m+1} a_{i+m-1}$.
  \item $\sum_{i=m}^{n} a_{i}=\sum_{i=1}^{n} a_{i}-\sum_{i=1}^{m-1} a_{i}$.
  \item Pour $c \in \mathbb{R}, \sum_{i=m}^{n} c=(n-m+1) c$.
  \item Pour $c \in \mathbb{R}, \sum_{i=m}^{n} c a_{i}=c \sum_{i=m}^{n} a_{i}$.
  \item $\sum_{i=1}^{n}\left(a_{i}+b_{i}\right)=\left(\sum_{i=1}^{n} a_{i}\right)+\left(\sum_{i=1}^{n} b_{i}\right)$
  \item $\sum_{i=1}^{n}\left(a_{i}-a_{i-1}\right)=\answer{a_{n}-a_{0}}$. telescoping sum
  \item $\sum_{i=1}^{n} i=\answer{\frac{n(n+1)}{2}}$.
  \item $\sum_{i=1}^{n} i^{2}=\frac{n(n+1)(2 n+1)}{6}$.
  \item $\sum_{i=1}^{n} i^{3}=\frac{n^{2}(n+1)^{2}}{4}$.
  \item $\sum_{i=1}^{n} r^{i}=\frac{r^{n+1}-r}{r-1}$.
\end{enumerate}
\end{proposition}

\begin{example}\label{example:12-5}
% work: sum_{k=5}^{12} k^2 = sum_{k=1}^{12} k^2 - sum_{k=1}^{4} k^2 = 12*13*25/6 - 4*5*9/6 = 650 - 30 = 620
\[
\begin{aligned}
\sum_{i=1}^{100} i & =\frac{100 *(100+1)}{2}=\frac{10100}{2}=5050 \\
\sum_{k=5}^{12} k^{2} & = \answer{620}
\end{aligned}
\]
\end{example}

\begin{exercise}\label{exercise:12-6}
Try one with a partner.
% work: sum (i^2 - i - 10) for i=1..20 = 20*21*41/6 - 20*21/2 - 10*20 = 2870 - 210 - 200 = 2460
\[
\sum_{i=1}^{20} i^{2}-i-10 = \answer{2460}
\]
\end{exercise}

\end{document}
